\subsection{Saturation-index stratification and nongeneric savings}
\label{sec:flag-defect-stratification}

Keep the local notation of Section~\ref{sec:flag-generic-local-amplitude},
including \(s=r+1\), \(m=2s\), \(n=k+2r+1\), \(l=n-m=k-1\),
and \(P=\mathrm{Hp}=q^{v(\theta)+a+d}\). We work with the actual
admitted term lattices and their two stabilizer weights.
Assume \(k\ge4r+7\), and that the point has a term.
For a fixed frame-class representative let
\[
 M=M(u,u'),\qquad M_h=\Lambda\cap J,\qquad
 \beta=[M_h:M]=q^i.
\]
All indices here are additive cardinalities, in particular at a split
place. By (GL3), \(0\le i\le a\). A term is nongeneric precisely when
\(i\ge1\). We establish a gain without any phase estimate on the
strict-image branch of the defect stratification, and on the
large-defect branch of the remaining terms.

\begin{lemma}[The complement image]
\label{lem:flag-defect-image}
Put \(W=J^\perp\), and let \(L_\delta\subset W\) be the
intersection of the two inner quotient lattices, with the additional
condition that the pair of quotient classes \((w,w)\) is in the
prescribed level graph when \(\theta\ne1\).
For a term over \(\delta\) this is a lattice containing
\(L=\Lambda\cap W\), and there is a canonical \(O\)-linear injection
\[
 L_\delta/L\ \longrightarrow\ M_h/M,\qquad
 w\longmapsto [\theta^{-1}(x'-x)],
 \tag{DS1}
\]
where \(x\in X,\ x'\in X'\) satisfy
\(w+x,w+x'\in\Lambda,\ x'-x\in\theta\Lambda\).
Its image is the original \(\mathrm{Ihat}/M\), so
\([L_\delta:L]=\mathrm{cardI}=q^j\), with \(0\le j\le i\).
For a nongeneric term which is not of the original
\(\mathrm{TypeBe}\), one has \(j<i\).
\end{lemma}
\begin{proof}
The quotient maps embed \(W\) in each inner quotient. Their
inverse images of the two inner lattices are full lattices there;
the graph condition is a linear congruence. Equivalently
\(L_\delta\) consists of the \(w\) admitting the indicated lifts:
at ordinary places this is just the two membership conditions
because \(\theta=1\). This is exactly the source's complement
datum and \(\mathrm{TauRel}\).

The difference of two choices of \(x\) belongs to
\(\Lambda\cap X=\sum_{b<s-1}Ou_b\), and similarly on the
primed side, by primitivity of the two adapted frames.
Write the two differences as \(u c\) and \(u' c'\), using
only the first \(r\) entries. If \(\theta=1\) their difference
is already in \(M\). At a level place the lift condition
gives \(u'c'-uc\in\theta\Lambda\). Since \(u'\equiv u\)
modulo \(\theta\Lambda\) and \(u\) is primitive,
\(c'\equiv c\bmod\theta O^r\). Thus
\[
 \theta^{-1}(u'c'-uc)=u\,\theta^{-1}(c'-c)
                          +\theta^{-1}(u'-u)c'\in M .
\]
This proves independence of the lifts. Add and scale lifts to
prove linearity. The value belongs to \(M_h/M\).
If it is zero, the independence of the two joint frames
shows \(x,x'\) are integral combinations of their own
first \(r\) vectors: indeed the coefficients in
\(\theta M\) are of the form \((\theta c-b,b)\)
with \(c,b\in O^s\). Hence \(w\in L\).
Conversely \(w\in L\) admits zero lifts. This proves (DS1)
and the asserted image description, since the image is
already an \(O\)-module.

For the last assertion, suppose \(i>0\) and \(j=i\).
Every element of \(M_h/M\) then has a representative
\(\theta^{-1}(x'-x)\) from the two isotropic spans.
The two last-vector coefficients of \(\theta y\) for
any \(y\in M_h\) therefore satisfy the source's
\(\mathrm{RR}\) condition: the primed coefficient
is integral and the sum of the two coefficients belongs
to \(\theta O\). This proves \(\mathrm{levC}(y)\le r\).
For the other condition in \(\mathrm{TypeBe}\), take
\(a_\nu\in X,\ a'_\nu\in X'\) with
\(a_\nu+a'_\nu\in\theta\Lambda\), for \(\nu=1,2\).
By surjectivity of (DS1) choose \(w\in L_\delta\) with
lifts \(x,x'\) such that
\[
 (a_1+a'_1)-(x'-x)\in\theta M .
\]
Comparison of the independent frame coefficients shows
\(a_1+x\in\sum_{b<r}Ou_b\subset\Lambda\).
Also \(w+x\in\Lambda\), with \(w\perp J\), so
\(\langle x,a_2+a'_2\rangle\in\theta O\).
Because \(X\) is totally isotropic this implies
\(\langle x,a'_2\rangle\in\theta O\).
Similarly \(\langle a_1+x,a'_2\rangle\in\theta O\).
Subtracting proves \(\langle a_1,a'_2\rangle\in\theta O\).
Together with nongenericity and the level condition this is
exactly \(\mathrm{TypeBe}\), a contradiction.
\end{proof}

\begin{lemma}[Uniform sublattice count]
\label{lem:flag-defect-sublattices}
For a free rank-\(l\) lattice over \(O\), the number of
sublattices of additive index \(q^j\) is at most
\[
             16(j+1)q^{(l-1)j}. \tag{DS2}
\]
In the unramified nonsplit case it is zero for odd \(j\),
and at most \(4q^{(l-1)j}\) otherwise.
\end{lemma}
\begin{proof}
Over a discrete valuation ring of residue size \(Q\),
successive elimination in a fixed ordered basis gives
the triangular normal forms with diagonal valuations
\(d_1,\dots,d_l\ge0\), \(\sum d_\nu=d\), and, for fixed
valuations, exactly \(Q^{\sum_{\nu=1}^l(\nu-1)d_\nu}\)
choices of off-diagonal residue entries (reverse the
basis order if using the other triangular convention).
The number of index \(Q^d\) sublattices is therefore
at most
\[
 Q^{(l-1)d}\prod_{\nu=1}^{l-1}(1-Q^{-\nu})^{-1}
       <4Q^{(l-1)d}.
\]
For the last bound use \(Q\ge2\) and
\(\prod_{\nu\ge3}(1-2^{-\nu})\ge
1-\sum_{\nu\ge3}2^{-\nu}=3/4\);
the first two factors multiply to \(3/8\).
At an inert place \(Q=q^2\); at a split place distribute
\(j\) as \(j_1+j_2\), count each factor with \(Q=q\),
and sum over the \(j+1\) possible decompositions.
This proves (DS2).
\end{proof}

\begin{theorem}[Strict-image and large-defect amplitude gain]
\label{thm:flag-defect-stratification}
Write \(A_\delta(\mathcal P)\) for the original phase-weighted
amplitude restricted to a predicate \(\mathcal P\) on the
actual terms, and
\[
 \|\mathcal P\|_{\rm amp}^2
       =\sum_\delta \mathrm{wt}(\delta)|A_\delta(\mathcal P)|^2 .
\]
The following bounds hold in the original local rank range:
\[
 \begin{split}
 \|\,i\ge1,\ j<i\,\|_{\rm amp}^2
    &\le P^{2s^2+10s+11}q^{-l}\,nn,\\
 \|\,i\ge I\,\|_{\rm amp}^2
    &\le P^{2s^2+10s+11}q^{-I}\,nn
                      \qquad(I\ge1).
 \end{split}                                                   \tag{DS3}
\]
Both also hold for any further invariant restriction of the
indicated terms. The union of the strict-image branch and
the large-defect branch \(i\ge k-2r-5\) satisfies the original
nongeneric amplitude bound
\[
 \|\,i>j\ \text{or}\ i\ge k-2r-5\,\|_{\rm amp}^2
        \le P^{39s^2}q^{-(k-2r-5)}\,nn.                    \tag{DS4}
\]
When \(P=1\) all these predicates are empty. The only
nongeneric terms not covered by (DS4) have
\(1\le i=j<k-2r-5\); they belong to \(\mathrm{TypeBe}\).
\end{theorem}
\begin{proof}
Set \(e=s^2+3s\), \(e_0=s^2+7s\).
For a fixed integral frame class \(C\) the possible
\(M_h\) are at most \(q^{2am}\) by (GL3).
For fixed \(\delta,M_h\) and \(j\), the complement
\(L\) is a sublattice of the fixed \(L_\delta\)
of index \(q^j\). Lemma~\ref{lem:flag-defect-sublattices}
bounds the number of choices. For fixed \(M_h,L\),
the anti-isometric discriminant gluing argument of
(GL5), now applied to \(M_h\), bounds the number of
self-dual lattices by \(q^{2am}\): its discriminant
module is a quotient of a submodule of \(M^\#/M\)
and is generated by at most \(m\) elements over \(O\).
Thus (GL4), followed by the exact orbit weights
of (GL6), gives, for fixed \(\delta\) and \(i,j\),
\[
 m_\delta(i,j)\le
 16(j+1)P^{e_0}q^{(l-1)j},
 \qquad e_0=s(s-1)+4m.
 \tag{DS5}
\]
Indeed \(q^{4am}\le P^{4m}\), so this bound has no
hidden dependence on the ambient rank. All counts
are upper bounds, permitting every further restriction.

Independently, summing over the complement orbits by
(GL6) and applying (GL1) to the overlattices of index
\(q^i\) gives
\[
 \sum_\delta \mathrm{wt}(\delta)m_\delta(i)
       \le \tfrac54 c_x P^e q^{-il}
       \le \tfrac53 P^e q^{-il}nn.                  \tag{DS6}
\]
Here \(c_x\) is the common Haar scalar for all corrected
frame classes. For the last step use any occupied
reference class \(C_*\):
\(nn\ge ngen(C_*)=w(C_*){\cal M}(M(C_*))
\ge3c_x/4\) by (GL6), (GL4) and (GL1).
This does not require that the reference term itself
be generic.

The modulus of an amplitude over a subset is at most
its mass. For \(i\ge1\), sum (DS5) over \(0\le j<i\)
and use \(\sum_{h=0}^{b} q^{(l-1)h}\le
2q^{(l-1)b}\) (here \(l\ge6\)).
The pointwise mass bound is
\(32 i P^{e_0}q^{(l-1)(i-1)}\).
Multiplying by (DS6) bounds the squared amplitude norm
on this \(i\)-stratum by
\[
 \tfrac{160}{3} i P^{e_0+e}q^{-l-(i-1)}nn .
\]
For the whole stratum \(i\ge1\), summation over \(j\le i\)
instead gives the bound
\(\tfrac{160}{3}(i+1)P^{e_0+e}q^{-i}nn\).
These are inequalities for the original two
reciprocal-stabilizer weights; they do not count
unweighted representatives of those groups.

Use the triangle inequality in weighted \(\ell^2\)
over the data for the sum of amplitudes at the distinct
indices \(i\), not the sum of their absolute squares.
Since \(i\le a\),
\[
 \sum_{i\ge I}^{a}\sqrt{i+1}\,q^{-i/2}
 \le\frac{\sqrt{a+1}\,q^{-I/2}}{1-2^{-1/2}},
\]
and the analogous strict-image sum has the same bound
with \(q^{-I/2}\) replaced by \(1\), after removing
\(q^{-l/2}\). The square of the denominator inverse
is \(6+4\sqrt2<12\). Consequently the squared norms
are bounded by \(640(a+1)P^{e_0+e}nn\) times
the respective gains. In a nonempty nongeneric branch
\(P\ge q^a\ge2\), \(a+1\le P\), and \(640\le 2^{10}\).
This proves (DS3), with \(e_0+e=2s^2+10s\).
For (DS4), disjointly assign strict-image terms to
the first estimate and the remaining terms with
\(i\ge k-2r-5\) to the second. Cauchy--Schwarz for
the sum of these two amplitudes costs at most \(4\le P^2\).
The rank inequality gives \(l\ge k-2r-5\), and
\(2s^2+10s+13\le39s^2\) for every \(s\ge1\).
Thus (DS4) follows. Lemma~\ref{lem:flag-defect-image}
identifies the missing terms.
\end{proof}

For example, taking \(I=1\) in (DS3) already gives a
uniform nongeneric gain \(q^{-1}\), without cancellation
and without enlarging the original height exponent.
At an inert place the same proof starts at \(I=2\).
For \(a<k-2r-5\) the large-defect branch is empty;
this is not a proof that the remaining terms vanish.
In that case the unresolved local problem is precisely
phase cancellation on the small full-image defect strata.
The source's second-type family estimate
\(\mathrm{amplitude\_le\_mass\_secondType\_family}\)
supplies the stronger gain on that branch, but its
character calculation has not been reconstructed here.
